\chapter{Standardizing Explanation Regions for Masking-Based Tests (Day 4)}
\label{chap:day04}

\section{Motivation}

Day 3 left every sample with a baseline judgment and, for rules~1--3, a person-relative detection
process that grounds independent worker and safety-object boxes rather than asking a single
scene-level existence question. Those boxes were produced to decide compliant/violation; they were
not yet organized into anything a later explanation-testing method could consume. Descriptive
Accuracy, the masking-based evaluation planned for Day 5, needs a single, well-defined region to
mask per sample --- a candidate answer to "which part of the image does the model's judgment depend
on?" --- and a procedure for actually producing the masked image. Day 4's task was to build that
standardization layer: take Day 3's already-computed boxes, rank them into an ordered list of
candidate explanation regions, and provide a masking primitive that blacks out (or blurs) the
top-ranked region for re-inference. Day 3 and Day 4's work were developed together, per commit
\texttt{6de76fb}, reflecting that the redesigned proxy and the regions layer that consumes its boxes
were designed as a single piece of work rather than as two independent days.

\section{Implementation}

\texttt{regions.py} introduces a \texttt{Region} dataclass --- a \texttt{box}, a \texttt{source}
field (\texttt{"model"} or \texttt{"grid"}), and a \texttt{label} --- as the standard unit the rest
of the pilot's masking-based tests operate on. \texttt{standardize\_regions(boxes, labels,
image\_size, grid=(4, 4))} is the layer's central function: it ranks model-returned boxes by
\textbf{area, descending}, and returns them as a list of \texttt{Region} objects with
\texttt{source="model"}. Florence-2's grounding output carries no attention or saliency map, so
there is no real per-region importance score available to rank by; area is used as an explicit,
acknowledged stand-in, on the assumption that a larger detection is more likely to be the salient
one than a sliver. Only when the model returned \emph{zero} boxes at all does
\texttt{standardize\_regions} fall back to an unranked $4\times4$ grid, via
\texttt{grid\_fallback\_regions}, so that every sample still has some candidate region to mask even
when grounding failed outright.

Three smaller helpers round out the module: \texttt{iou()} computes intersection-over-union between
two boxes; \texttt{mask\_region()} fills a given box with either a solid black patch or a
Gaussian-blurred copy of the same region (\texttt{mode="black"} or \texttt{mode="blur"}); and
\texttt{boxes\_overlap\_or\_close()} tests whether two boxes overlap or have edges within a distance
threshold. The last of these underlies \texttt{all\_boxes\_covered()}, which is, in fact, the
function that Chapter 3's redesigned \texttt{\_answer\_presence\_rule} actually calls to check that
every detected worker is covered by a nearby safety-object box. That helper's home in
\texttt{regions.py} rather than in Day 3's own module is the clearest evidence that Day 3's proxy
redesign and Day 4's regions layer were not sequential, independent pieces of work but a single
designed unit, split across two chapters here only for narrative pacing.

\texttt{scripts/04\_extract\_regions.py} applies \texttt{standardize\_regions} to all 163 samples,
reusing the \texttt{worker\_boxes} and \texttt{object\_boxes} already written to
\texttt{results/baseline\_predictions.csv} by Day 3 --- no new GPU inference was needed for this
step. For each sample it records the top-ranked region's source and label, and for the first ten
samples it additionally saves a before/after masked-image pair, applying \texttt{mask\_region} in
\texttt{"black"} mode to the top region, so the masking step itself can be sanity-checked visually
before Day 5 relies on it.

\section{Architecture: A New Masking Layer}

Figure~\ref{fig:diagram-c} shows where this chapter's layer sits relative to Day 3's baseline.
Day 3 produces grounded boxes and labels; Day 4 adds everything downstream of that point ---
ranking the boxes into an ordered region list, and masking the top-ranked region to produce an
image suitable for the re-inference step that Descriptive Accuracy will perform on Day 5. Nothing
in this diagram changes how Day 3's boxes are produced; it only standardizes what is done with them
afterward.

\begin{figure}[h]
\centering
\resizebox{\textwidth}{!}{%
\begin{tikzpicture}[node distance=1.6cm]
\node[block] (boxes) {Baseline boxes\\+ labels (Day 3)};
\node[block, right=of boxes] (rank) {standardize\_regions\\(rank by area /\\grid fallback)};
\node[block, right=of rank] (region) {Ranked\\Region list};
\node[block, right=of region] (mask) {mask\_region\\(black / blur)};
\node[block, right=of mask] (out) {Masked image\\for re-inference};
\draw[arrow] (boxes) -- (rank);
\draw[arrow] (rank) -- (region);
\draw[arrow] (region) -- (mask);
\draw[arrow] (mask) -- (out);
\end{tikzpicture}%
}
\caption{Regions/masking layer added on top of Day 3's baseline grounding output.}
\label{fig:diagram-c}
\end{figure}

\section{Worked Example}

Figure~\ref{fig:day04-worked} revisits image 7's rule~1 (hard hat) sample, the same worked example
used in Chapter 3, for narrative continuity. The left panel shows the top-ranked region --- the
largest of the model-returned boxes, outlined for inspection --- overlaid on the original image; the
right panel shows the result of \texttt{mask\_region} applied to that same box in
\texttt{"black"} mode. The masked region corresponds exactly to the ranked top box and nothing else
in the frame, confirming that the masking primitive is doing precisely what it is meant to do before
it is relied on for the Descriptive Accuracy measurements in Chapter 5.

\begin{figure}[h]
\centering
\begin{subcaptionblock}{0.48\linewidth}
\centering
\includegraphics[width=\linewidth]{regions/0000007_rule_1_before.png}
\caption{Before: top-ranked region outlined.}
\end{subcaptionblock}
\hfill
\begin{subcaptionblock}{0.48\linewidth}
\centering
\includegraphics[width=\linewidth]{regions/0000007_rule_1_after.png}
\caption{After: top-ranked region masked.}
\end{subcaptionblock}
\caption{Image 7, rule~1 (hard hat): the region ranked first by area, before and after
\texttt{mask\_region} blacks it out.}
\label{fig:day04-worked}
\end{figure}

\section{Results}

Running \texttt{standardize\_regions} over all 163 samples shows that a real model-returned box was
available and ranked first for \textbf{159/163 samples (97.5\%)}; the remaining \textbf{4/163
(2.5\%)} returned zero boxes and correctly fell back to the unranked grid. The model-box rate this
high confirms that Day 3's grounding calls are, for the overwhelming majority of samples, supplying
the regions layer with something real to rank, and that the grid fallback is functioning as the
rare exception it was designed to be rather than as a frequently-used default. Visual inspection of
2 of the 10 saved before/after pairs confirmed that \texttt{mask\_region} blacks out exactly the
ranked top box and nothing else in the frame, as shown in Figure~\ref{fig:day04-worked}.

\section{Problems Encountered and Resolutions}

The central design problem on Day 4 was not a bug but a missing capability in the underlying model:
Florence-2's grounding output provides no real saliency signal, no per-pixel or per-region
importance score of the kind that an attention map or a gradient-based attribution method would
supply. Without such a signal, there was no principled way to rank multiple candidate boxes by
"how much the model's judgment depends on this region." The resolution was an explicit design
decision, made and documented here rather than left implicit, to rank by box area as the best
available stand-in --- on the working assumption that a larger detected region is more likely to be
the one driving the judgment than a small one --- while flagging plainly that this is a stand-in,
not a claim of true importance. \textbf{This choice is shown later in the pilot (Chapters 6, 7, and
10) to systematically disadvantage small PPE objects}: a hard hat or harness is small relative to
the worker or scene around it, so an area-based ranking will tend to rank larger, non-PPE boxes
above the actual PPE object whenever both are present as candidates. That consequence is
forward-referenced here, at the point the choice was made, so it reads in this report as a foreseen
limitation rather than one discovered by accident downstream.

\section{Standard vs. Non-Standard Choices}

\textbf{Standard}: using bounding-box masking or occlusion as an explanation-testing primitive is a
long-established pattern in vision explainability, the same basic idea behind occlusion sensitivity
maps --- block out a region, observe whether and how the model's output changes, and treat a change
as evidence that the region mattered. Adopting that primitive here, rather than inventing a new
testing mechanism, follows established practice.

\textbf{Non-standard}: ranking candidate regions purely by box area as a substitute for a true
importance score is not a published method, and is not presented as one. It is an explicit,
acknowledged simplification, adopted because Florence-2's grounding output leaves no real
saliency signal to rank by, and reported here as a designed-and-flagged limitation rather than
folded silently into the pipeline as if it were equivalent to a genuine importance ranking.

\section{Implications for the Paper}

The area-ranking heuristic should be flagged explicitly in the paper as a limitation and
future-work item: it is the single highest-leverage fix identified across the entire pilot (see
Chapter 10, Finding 1, and Chapter 14's task list, item 3). Its effect is not confined to this
chapter. Because \texttt{standardize\_regions}'s ranking feeds every later masking-based test built
on top of it, the area-based choice made here is the mechanism behind the PPE-disadvantage findings
reported in Chapters 6, 7, and 10: small, genuinely safety-relevant objects are structurally
under-ranked relative to larger, often less relevant boxes, independent of how well any later
evaluation metric itself is designed. Recording that link here, at the point of the original design
decision, lets the paper present the limitation as anticipated and explained rather than as a
finding that surfaces without an account of its cause.
